% ============================================================================
%  SECCIÓN 1.8: JUSTIFICACIÓN
% ============================================================================

\section{Justificaci\'on}

La realizaci\'on de este proyecto se justifica por la importancia de acercar la
infraestructura deportiva del pa\'is a la poblaci\'on mediante el uso de
tecnolog\'ias de informaci\'on accesibles, masivas y de bajo costo. Justificar es
exponer los motivos que ameritan realizar la investigaci\'on, respondiendo a
las preguntas de por qu\'e es importante su realizaci\'on y para qu\'e se realiza
\citep{arias2012}. En ese sentido, el proyecto aporta desde las dimensiones
social, econ\'omica, t\'ecnica, cient\'ifica y metodol\'ogica.

\subsection{Justificaci\'on social}

El proyecto beneficiar\'a directamente a la ciudadan\'ia al facilitar el acceso a
la informaci\'on sobre los escenarios deportivos, promoviendo la pr\'actica del
deporte y la actividad f\'isica como factores de salud, inclusi\'on y convivencia
social. Al visibilizar la oferta deportiva existente, se contribuye a reducir
las barreras de acceso y a fomentar h\'abitos de vida saludable en la
poblaci\'on, en particular en j\'ovenes y ni\~nos. Asimismo, se favorece la
participaci\'on de la comunidad en los eventos deportivos que se organizan en
cada regi\'on, fortaleciendo el tejido social.

\subsection{Justificaci\'on econ\'omica}

Desde el punto de vista econ\'omico, el proyecto permite optimizar el uso de la
infraestructura deportiva existente, cuya subutilizaci\'on representa una
p\'erdida de recursos para las entidades administradoras. Al dar visibilidad a
los escenarios, se incrementan las posibilidades de uso, alquiler y
organizaci\'on de eventos, generando ingresos para los administradores y
dinamizando la econom\'ia local vinculada al deporte. Adem\'as, el desarrollo de
una soluci\'on multiplataforma con herramientas de c\'odigo abierto y servicios
en la nube de bajo costo hace viable el proyecto con una inversi\'on m\'inima.

\subsection{Justificaci\'on t\'ecnica}

El proyecto aprovecha tecnolog\'ias maduras y ampliamente difundidas para el
desarrollo de aplicaciones web y m\'oviles, sistemas de informaci\'on geogr\'afica
y mapas interactivos \citep{longley2015, pressman2010}. Estas tecnolog\'ias han
sido validadas en soluciones similares de otros sectores y permiten construir
una aplicaci\'on multiplataforma (iOS y Android) con un \'unico c\'odigo base,
reduciendo costos y tiempos de desarrollo. La soluci\'on propuesta representa
una innovaci\'on en el \'ambito deportivo nacional, donde no existe actualmente un
sistema equivalente de acceso p\'ublico.

\subsection{Justificaci\'on cient\'ifica}

Los argumentos cient\'ificos que sustentan la realizaci\'on del proyecto se
fundamentan en el cuerpo te\'orico de la ingenier\'ia de software, los sistemas
de informaci\'on geogr\'afica y la metodolog\'ia de la investigaci\'on
\citep{sommerville2011, longley2015, hernandez2014}. El proyecto aplica
principios establecidos de an\'alisis y dise\~no de sistemas, usabilidad y
calidad de software, y genera nuevo conocimiento sobre la aplicaci\'on de estas
herramientas al \'ambito de la difusi\'on de la infraestructura deportiva,
obteniendo procedimientos que permiten un acercamiento eficaz al objeto de
estudio.

\subsection{Justificaci\'on metodol\'ogica}

El proyecto seguir\'a un enfoque de desarrollo \'agil e iterativo, acorde con las
buenas pr\'acticas de la ingenier\'ia de software \citep{sommerville2011,
larman2004}. La definici\'on del producto m\'inimo viable permite validar la
propuesta con usuarios reales en etapas tempranas y mejorar de forma
continua, aplicando metodolog\'ias que podr\'an servir de referencia para
proyectos similares de digitalizaci\'on de la informaci\'on deportiva.

\subsection{An\'alisis de viabilidad}

De acuerdo con la gu\'ia del avance, se analiza la viabilidad del proyecto
considerando el tiempo disponible, la cantidad de integrantes, los
conocimientos actuales, las tecnolog\'ias disponibles y la complejidad de la
soluci\'on:

\begin{itemize}
  \item Tiempo: el proyecto se desarrollar\'a durante la gesti\'on
    acad\'emica, con un cronograma de 16 semanas distribuidas en cinco fases
    (an\'alisis, dise\'no, desarrollo, pruebas y cierre).
  \item Recursos humanos: el equipo est\'a conformado por un m\'aximo de
    seis integrantes con roles definidos (l\'ider, analista, dise\'nador UI/UX,
    desarrollador y gestor de documentaci\'on).
  \item Conocimientos: las tecnolog\'ias propuestas (React Native,
    Expo, Node.js, PostgreSQL y Leaflet) se alinean con los contenidos de la
    asignatura y con las competencias del equipo.
  \item Tecnolog\'ias disponibles: se emplean herramientas de c\'odigo
    abierto y servicios en la nube con planes gratuitos, sin requerir
    inversi\'on econ\'omica significativa.
  \item Complejidad: el alcance se limita a dos o tres perfiles de
    usuario y entre cinco y ocho funcionalidades principales, definidas en el
    producto m\'inimo viable.
\end{itemize}

En consecuencia, el proyecto es viable: el problema est\'a claramente
definido, la soluci\'on aporta valor a los usuarios, el n\'umero de usuarios y
funcionalidades es manejable, y el MVP puede desarrollarse y demostrarse de
forma funcional durante el m\'odulo.
